%% kapa: exam=1 round=31 year=2020 period=2 form=A subject=accounting subject_ko=회계학 seq=2 number=42
%% source: Q-Net 감정평가사 기출문제 (cst00309) artlSeq=5207219 file=extracted/1008/5207219/A형/A형/2020년도 제31회 감정평가사 1차 시험_2교시_A형.pdf
%% conversion: kordoc 4.24.51 -> tools/gen1.py (auto); review=proofread-sonnet; flags=none
\begin{question}{42}
재무제표 표시에 관한 설명으로 옳은 것은?
\begin{choices}
  \item 비용을 성격별로 분류하는 경우에는 적어도 매출원가를 다른 비용과 분리하여 공시해야 한다.
  \item 기타포괄손익의 항목(재분류조정 포함)과 관련한 법인세비용 금액은 포괄손익계산서에 직접 표시해야 하며 주석을 통한 공시는 허용하지 않는다.
  \item 유동자산과 비유동자산을 구분하여 표시하는 경우라면 이연법인세자산을 유동자산으로 분류할 수 있다.
  \item 한국채택국제회계기준에서 별도로 허용하지 않는 한, 중요하지 않은 항목이라도 유사항목과 통합하여 표시해서는 안된다.
  \item 경영진은 재무제표를 작성할 때 계속기업으로서의 존속가능성을 평가해야 한다.
\end{choices}
\end{question}
\answer{5}
